\section{Conclusion}\label{sec:conclusion}

We have answered, affirmatively and constructively, whether alignment dynamics admit a
variational free-energy model in which a narrow misalignment intervention induces a phase
transition. Modeling alignment as an FEP gradient flow (Theorem~\ref{thm:gradflow}) and deriving
its potential from a generative mixture (Proposition~\ref{prop:micro}), we showed that the
universal reduction is the cusp catastrophe \eqref{eq:cusp}. Within this single family we proved
three things. A unique critical threshold $\lambda^\ast$ provably separates robust alignment from
emergent misalignment, located exactly at a Hessian eigenvalue crossing
(Theorem~\ref{thm:threshold}). The order of the transition is set by the symmetry of the
intervention: continuous for symmetric interventions (Theorem~\ref{thm:pitchfork}) and a
catastrophic jump with hysteresis for biased ones (Theorem~\ref{thm:fold}), unifying the
gradual-LoRA and sudden-full-finetune regimes in one normal form. And near the threshold the
variance and autocorrelation time of the order parameter diverge
(Theorem~\ref{thm:earlywarning}), giving a critical-slowing-down early-warning law in which the
internal order parameter leads the behavioral failure.

The key takeaway is a dictionary, backed by theorems and by symbolic and numerical verification,
between the Free Energy Principle, critical-transition theory, and emergent misalignment: it
turns an empirical surprise into an a priori, monitorable risk calculation, with null-versus-EM
classification at ROC-AUC $0.996$ and threshold MAE $0.056$.

Future work should estimate the landscape parameters $(a_0,c,h)$ directly from model internals---
from basin curvature and from the alignment of the finetuning direction with the misalignment
axis---to make the threshold prediction operational on real training runs, and should test
whether the misalignment-onset universality class is genuinely mean-field or carries the
nontrivial exponents seen in related learning transitions.
